\documentclass[11pt]{article}
\usepackage[utf8]{inputenc}
\usepackage[T1]{fontenc}
\usepackage{lmodern}
\usepackage[margin=1in]{geometry}
\usepackage{csquotes}
\usepackage[style=authoryear,backend=biber,sorting=nyt,maxbibnames=99,giveninits=true,uniquename=init]{biblatex}
\usepackage[hidelinks]{hyperref}
\addbibresource{references.bib}
\begin{document}
\section*{Reactive Media Nutrient Removal}

\subsection{Overview}

Reactive-media treatment technologies intercept nutrient loads in subsurface (tile) drainage before they reach surface waters, using carbon-rich media (woodchips) to drive microbial denitrification of nitrate and iron- or calcium-rich media to adsorb or precipitate dissolved phosphate. Two decades of controlled microcosm and column studies, pilot- and full-scale field installations, and global syntheses have established these practices---especially denitrifying woodchip bioreactors---among the best-characterised edge-of-field measures for artificially drained cropland. The consolidated picture is that woodchip bioreactors remove nitrate dependably across a wide range of climates and designs, while reactive filter media retain dissolved reactive phosphorus (DRP) that would otherwise pass through drains in solution. Consolidating full-size systems worldwide, \textcite{Christianson2021} report a mean nitrate-N load reduction of $40 \pm 26\%$ (median 46\%; $n=15$) for bed-style bioreactors and removal rates of $7.2 \pm 9.6$ g N m$^{-3}$ d$^{-1}$ (median 5.1; $n=27$), with the practice assessed as removing 20--40\% of annual nitrate-N loss in the Midwest. The largest compilation to date---36 site-years at ten full-size Illinois bioreactors---found a median annual removal of 49 kg N per bioreactor per year (4.1 kg N ha$^{-1}$ y$^{-1}$), a median edge-of-field efficiency of 17\%, and removal in \emph{all} years, confirming dependability while pointing to room for optimisation \parencite{Christianson2026}. For phosphorus, laboratory batch and column tests routinely achieve $\geq 90$--$99\%$ removal for the better media, industrial by-products adsorb roughly an order of magnitude more phosphate than natural minerals \parencite{Sellner2019b}, and field structures treating real flow have removed on the order of 27--55\% of the dissolved-P load \parencite{Penn2020,Penn2024,Shedekar2020}.

Several consistent lines of evidence organise what is known. Field monitoring quantifies concentration reductions, load reductions, and volumetric removal rates and their dependence on hydraulic residence time (HRT), temperature, flow, and loading. Mechanistic and microbial studies establish denitrification as the dominant nitrate pathway, sustained for years by woodchip carbon, and inner-sphere adsorption versus Ca-phosphate precipitation as the two phosphate pathways. An active design literature refines sizing, hydraulics, and enhancement strategies. Work on secondary effects documents by-products---nitrous oxide, methane, sulfate reduction, dissolved organic carbon (DOC), and phosphorus dynamics---and generally concludes that these ``pollution-swapping'' risks are real but small at landscape scale. A recurring qualifier is that further study is warranted for phosphorus, pesticides, and greenhouse gases \parencite{Mitchell2023}.

\subsection{Woodchip bioreactors: field-scale nitrate removal performance}

Field installations consistently reduce nitrate, but the magnitude spans a wide, well-explained range. In a controlled Illinois field experiment, bioreactors reduced nitrate loads by an average of 63\%, from 20\% at 2-h HRT to 98\% at 8-h HRT, with an average removal rate of 11.6 g NO$_3$-N m$^{-3}$ d$^{-1}$ (range 5--30) \parencite{Bell2015}. Under cold-climate conditions in Quebec, a three-year evaluation reduced effluent nitrate-N by 99\% and total-N by 72--73\% relative to influent, achieved a 99\% reduction in flow-weighted nitrate-N load, and reported median removal rates of 6.84 g NO$_3$-N m$^{-3}$ d$^{-1}$ (by volume) and 2.99 g NO$_3$-N m$^{-2}$ d$^{-1}$ (by area); nitrate-N was 84\% of influent total-N, non-compliance events fell by 96\%, and roughly 55\% of drainage passed through the bioreactors annually \parencite{Husk2017}. Across four South Dakota bioreactors, concentration reductions ranged from 7\% to 100\%, removal rates from 5.2 to 26.5 g N m$^{-3}$ d$^{-1}$, and load reductions from 39\% to 89\% \parencite{Thapa2023}. A New Zealand system raised removal efficiency from 31\% at 2.8-d HRT to 93\% at 6.4-d HRT \parencite{Rivas2020}, and a Missouri bioreactor typically produced effluent below 10 mg NO$_3$-N L$^{-1}$, cutting spring influent means of 17 mg L$^{-1}$ to 5 mg L$^{-1}$ and post-fertilisation means of 69 mg L$^{-1}$ to 21 mg L$^{-1}$ \parencite{Aide2016}.

Long-term and multi-site datasets place these numbers in context. In the Illinois ten-bioreactor compilation, systems removed a median 51\% of the N routed into them at a median annual rate of 4.6 g N m$^{-3}$ d$^{-1}$; standardised designs achieved a median flow-weighted HRT of 2.8 h (near the 3.0-h target), and only annual N-loss removal differed significantly between standardised and hydraulic-maximised designs (4.7 versus 3.2 kg N ha$^{-1}$; $p=0.031$) \parencite{Christianson2026}. An early nine-year Iowa denitrification-wall study found annual nitrate loss of 24.5 kg NO$_3$-N ha$^{-1}$ under the wall versus 54.5 kg ha$^{-1}$ under conventional drainage, with no significant trend over time \parencite{Moorman2010}. Not all systems meet targets: a Virginia bioreactor at the low end of N inputs cumulatively removed only $\sim$9.5\% of influent NO$_3$-N ($0.56 \pm 0.25$ g m$^{-3}$ d$^{-1}$) in its second year---greater in the growing (20.4\%) than non-growing (5.4\%) season---below the NRCS 25--45\% target, with net export on 17\% of paired-sample days \parencite{Bock2018}. Similarly, a paired Illinois system removed only 6--28\% of edge-of-field nitrate (1.8--5.6 kg N ha$^{-1}$; 52--161 kg N field$^{-1}$ annually) and met the federal 20\% edge-of-field target in only one of three years \parencite{Maxwell2022}. Controlled column studies confirm the achievable ceiling: complete removal at low flow (2.9 cm d$^{-1}$) declining to 30\% at 13.6 cm d$^{-1}$ \parencite{Greenan2009}, and 100\% removal over 100 days at 20 mg N L$^{-1}$ and 24-h HRT (10.1 g N m$^{-3}$ d$^{-1}$) \parencite{Hua2016}.

\subsection{Controls on bioreactor performance: hydraulics, temperature, and loading}

The dominant, repeatedly quantified control is hydraulic residence time, which trades off against the fraction of flow captured. Removal efficiency rises with HRT across studies---from 20\% (2 h) to 98\% (8 h) in Illinois \parencite{Bell2015}, from 31\% to 93\% over 2.8--6.4 d in New Zealand \parencite{Rivas2020}, and increasing at longer daily HRT but trending negative below about 2 h in a paired Illinois system \parencite{Maxwell2022}. In the ten-bioreactor dataset, annual flow-weighted HRT correlated strongly with flow captured ($\rho=-0.67$) and with treated-water N removal efficiency ($\rho=0.60$) but not with edge-of-field efficiency ($\rho=0.03$), as standardised-design theory predicts \parencite{Christianson2026}. Column work shows removal rate per unit wood increasing with flow while efficiency falls ($y=4.299x^{-1}$, $r^2=0.96$) \parencite{Greenan2009}, and heated-treatment tests showed removal rising from 1.5 mg NO$_3$-N L$^{-1}$ at 1 h to $>$18.6 mg L$^{-1}$ at 8 h \parencite{Ma2021}. Influent load also matters: higher loads raise removal \emph{rate} but lower percentage efficiency and outlet-concentration reductions---a pattern seen for bioreactors \parencite{Rivas2020,Christianson2026} and echoed in constructed-wetland syntheses \parencite{Mendes2021}. A meta-analysis of solid-carbon denitrification similarly found higher influent nitrogen associated with higher rates, and inferior removal above 5-h HRT for those studies \parencite{Zhang2021b}.

Temperature exerts a strong, consistent effect. In the only controlled Illinois field experiment holding influent nitrate and HRT constant, exponential functions of water temperature and HRT explained 73\% of the variance in load reduction and 43\% in removal rate \parencite{Bell2015}. In South Dakota, average removal was 56\% in March--June versus 81\% in July--August, falling by about 30\% below 12\,$^{\circ}$C \parencite{Thapa2023}; laboratory columns reached 96\% at 21.5\,$^{\circ}$C/24-h HRT versus 29\% at 10\,$^{\circ}$C/12-h HRT \parencite{Soupir2018}; and Michaelis--Menten $V_{max}$ dropped sharply from 22 to 5\,$^{\circ}$C for composted and fresh woodchips \parencite{Kouanda2021}. Cross-study analysis found low temperature ($<$20\,$^{\circ}$C) significantly depressed removal while intermediate and high temperatures did not differ \parencite{Zhang2021b}. Flow dynamics and antecedent conditions further modulate performance: removal was lower during high-flow events and at high inflow nitrate, and samples taken immediately after rainfall showed higher load removal than later samples \parencite{Thapa2023}. In Quebec, influent nitrate correlated with soil moisture ($r=0.64$) and prior-day temperature ($r=0.67$) and negatively with flow rate ($r=-0.53$), while in-bioreactor reduction tracked lagged temperature ($r=0.58$) and influent concentration ($r=0.71$) \parencite{Husk2017}. Reaction kinetics underpin these observations: removal follows zero-order kinetics when nitrate is non-limiting (rate constants 1.42--1.80 mg N L$^{-1}$ h$^{-1}$) and shifts to first-order (0.49 h$^{-1}$) once nitrate becomes limiting \parencite{Hua2016}, and can be treated operationally as zero-order when nitrate greatly exceeds $K_m$ \parencite{Kouanda2021}.

\subsection{Mechanisms, carbon substrates, and microbial ecology}

\subsubsection{Denitrification as the dominant pathway}

Denitrification is firmly established as the dominant nitrate-removal pathway. In Iowa microcosms it accounted for $>$98\% of transformed NO$_3$-N in all treatments (except cardboard, $>$96\%), with immobilisation $<$2.4\% and DNRA $<$1\% \parencite{Greenan2006}, and flow-through columns confirmed denitrification as the primary fate with immobilisation minor \parencite{Greenan2009}. Denitrification potential in woodchips is sustained and far exceeds surrounding soil: over nine years wood increased denitrification potential 31- to 4000-fold over subsoil, denitrifier populations exceeded $10^8$ g$^{-1}$ wood, and in situ push--pull testing gave nitrate half-lives of 0.64--0.81 d, implying 24-h HRT suffices to bring 20--25 mg L$^{-1}$ influent below 10 mg NO$_3$-N L$^{-1}$ \parencite{Moorman2010}. Reviews likewise identify denitrification as the main mechanism ahead of biological uptake and settling \parencite{Mendes2021}, while noting that DNRA (favoured under low NO$_3$-N and high carbon) and, more tentatively, anammox also occur \parencite{Hartfiel2022}.

\subsubsection{Carbon quality, substrates, and amendments}

Carbon quality and longevity govern long-term performance. Cellulose is the most rapidly consumed woodchip compound, making cellulose metabolism central to persistence \parencite{Hartfiel2022}. Woodchip decomposition is depth- and position-dependent: in an Iowa wall, wood loss reached 75\% at 90--100 cm (half-life 4.6 y) but $<$13\% at 155--170 cm (half-life 36.6 y) \parencite{Moorman2010}, and along a Minnesota bed the inlet showed greater decomposition (lower C/N, higher lignocellulose index) than downstream, with aged four-year woodchips more recalcitrant than fresh \parencite{Ghane2018}. Alternative and blended substrates can raise performance. Corncob-amended systems removed more nitrogen than woodchips alone: pilot CC75 reached 61.2\% at 16-h HRT versus 37.9\% for WC100, and modelled full-scale CC75 gave the highest annual load removal (15.6\%) and lowest cost, remaining cost-advantageous even with periodic corncob replacement \parencite{Law2023}. Corn stover blended up to 10\% with woodchips enhanced denitrification (mean concentration reductions of 86\% for corn stover, 89\% for the blend, 77\% for woodchips), though pure corn stover raised ammonium and total inorganic N and is not recommended as the main medium \parencite{Oladeji2023}. Microcosm rankings placed cornstalks $>$ cardboard $>$ wood plus soybean oil $>$ woodchips alone, with grinding and soybean-oil additions boosting removal \parencite{Greenan2006}, and Canadian-prairie testing found agricultural residues (barley straw best) outperforming woodchips for nitrate removal, though woodchips gave better hydraulic efficiency \parencite{Kohn2023}. Amendments further tune performance: activated carbon raised removal at low temperature ($<$10\,$^{\circ}$C) and cut organic-carbon losses while flaxseed cake did not improve removal despite greater carbon release \parencite{Povilaitis2020}; biochar increased nitrate removal in proportion to dose \parencite{Ahmadvand2020}; and composting woodchips before installation raised denitrification capacity 1.20--2.15-fold at 22\,$^{\circ}$C with better cold tolerance \parencite{Kouanda2021}. Broadly, synthetic solid carbon sources outperform natural ones, while mixed and single carbon sources did not differ significantly in meta-analysis, indicating species matters more than mixing \parencite{Zhang2021b}. A recurring minor by-product is nitrite accumulation (9.6--18.7\% of reduced nitrate), which increases as HRT shortens and influent nitrate rises \parencite{Hua2016,Kouanda2021}.

\subsubsection{Microbial ecology}

Bioreactor microbial communities are variable and shaped by substrate availability, depth, season, moisture, pH, temperature, inlet nitrate, dissolved oxygen, and HRT \parencite{Hartfiel2022}. Culture and molecular work has identified the responsible organisms, especially those active in cold conditions. From Minnesota bioreactors, screening of 206 isolates yielded three confirmed denitrifiers (\textit{Cellulomonas} sp.\ WB94, \textit{Delftia} sp.\ BE1.2, \textit{Microvirgula} sp.\ BE2.4), two of which denitrified aerobically and whose genomes carried denitrification and polysaccharide-degradation genes, marking WB94 and BE2.4 as bioaugmentation candidates \parencite{Anderson2020}. A parallel study enriched cold-active denitrifiers (\textit{Pseudomonas}, \textit{Polaromonas}, \textit{Cellulomonas}) and confirmed \textit{Cellulomonas} sp.\ WB94 reduces nitrate and degrades cellulose \parencite{Jang2019}. Nitrate-reducing fungi are also abundant and cold-adapted: densities reached 845{,}000 and 524{,}000 CFU g$^{-1}$ in soil and woodchips at 15\,$^{\circ}$C, 91 of 349 isolates produced gaseous N (frequently $^{29}$N$_2$, suggesting co-denitrification), 89\% showed cellulase activity, and \textit{nirK} and \textit{p450nor} were detected in 38\% and 12\% of strains \parencite{Aldossari2021}. Bioaugmentation shows conditional promise: pre-grown \textit{Pseudomonas} isolates from Danish woodchips increased nitrate removal by 27--37\% at 5\,$^{\circ}$C in a second assay, with \textit{P.\ chlororaphis} PSS4 most effective \parencite{Jeglot2022}, though inoculation alone often fails owing to rapid washout \parencite{Wang2025b}. Community function tends to be robust: under fluctuating water levels, denitrification rates and denitrifier communities were generally resistant to disturbance, and nitrous-oxide production was consistently about ten-fold lower than nitrate removal \parencite{Hathaway2017}. Mineral amendment can steer the community---heated tourmaline reduced biofilm growth yet increased the relative abundance of \textit{nirS}-, \textit{nirK}-, and \textit{nosZ}-encoding denitrifiers, reducing clogging while giving the highest removal rates (52.1, 43.1, 48.0 g NO$_3$-N m$^{-3}$ d$^{-1}$ at 1, 4, 8 h) \parencite{Ma2021}. Biofouling deposits in drainage wells were controlled chiefly by dissolved iron rather than aerobic bacteria or dissolved oxygen \parencite{Stuetz2004}.

\subsection{Bioreactor design, sizing, and enhancement strategies}

\subsubsection{Sizing and hydraulics}

Sizing procedures balance flow treatment against removal efficiency. Using 20\% of the Manning's-estimated peak flow generally treated the majority of annual volume (54\% at Hamilton, 72\% at SERF), and modelling supports design guidance of 10--20\% of peak flow at 6--8-h design retention times; for fixed volume, longer, narrower designs removed a greater percentage of load and had higher rates than short, wide designs \parencite{Christianson2013}. Load-duration analysis across three watersheds---where roughly 50\% of annual NO$_3$-N moved in the largest 30\% of daily base flows---showed that bioreactors sized to a 0.5-d HRT could remove about 44--52\% of base-flow nitrate and at least 20--30\% of total annual loss while occupying at most 0.27\% of watershed area \parencite{Moorman2015}. Treating the majority of annual flow is not universally necessary where goals are load-based: a similar design flow rate ($\sim$1.7--1.8 L s$^{-1}$) corresponded to 50\% non-exceedance in small drainage areas but only 5\% in large ones \parencite{Christianson2025}. Loading-density guidance from the Conservation Practice Standard was not supported by the ten-bioreactor data, which showed no significant relationship between loading density and removal efficiency \parencite{Christianson2026}.

Hydraulic behaviour links design to performance. Across 14 tracer tests at six full-size bioreactors, high-aspect-ratio conventional designs behaved most like plug flow (tanks-in-series), the wide design showed the most dispersion and short-circuiting, and an advanced wide design with baffles achieved higher volumetric efficiency and trended toward the highest removals; dispersion/short-circuiting and volumetric/hydraulic-efficiency metrics were each strongly correlated \parencite{deOliveira2023}. A wide baffled Illinois bioreactor nonetheless did not raise the fraction of flow treated ($\sim$40\%) or removal relative to conventional designs---attributed partly to flashy drainage (Richards--Baker index 0.74--0.91)---yet still achieved consistent 22--24\% edge-of-field reduction with 54--69.5\% internal efficiency and clear along-path nitrate declines \parencite{Dougherty2020b}. Expanding cross-sectional area improved flow capture: a paired short-and-wide system captured 84--92\% of annual discharge from a 29-ha field \parencite{Maxwell2022}, and wood-filled bioreactors showed better hydraulic efficiency than other feedstocks in the Canadian prairies \parencite{Kohn2023}. For buffer-diverting designs, more diverted flow did not mean more nitrate removal, buffer width reduced diversion nonlinearly, and neglecting exit head loss overestimated diverted flow and removal, motivating a process-based design with a removal component \parencite{Abdalaal2023}. Physical durability is quantifiable: hand-held LiDAR at 17 full-size bioreactors captured surface subsidence, which correlated with age (median 1.0 cm y$^{-1}$ at soil-covered versus 7.3 cm y$^{-1}$ at non-covered sites; maximum 73 cm) and with proximity to the drainage outlet \parencite{Christianson2025b}.

\subsubsection{Performance-enhancement strategies}

Several modifications extend or intensify treatment. Pumping can supply additional nitrate-laden water and lengthen the operating season: a retrofitted dual-chamber system pumped at 1.06--2.17 L s$^{-1}$ per chamber, extended operation by three weeks, and raised combined daily removal to 6.58 g N m$^{-3}$ d$^{-1}$ (from 3.62 without pumping), albeit bypassing 41\% of pumped flow untreated \parencite{Hartfiel2023}. At an underloaded Illinois bioreactor, pumping increased treated volume, loading, and removal---removing 2.8 kg N in a 29-d 2022 period (versus 0.69--2.0 kg N unpumped) and 7.2 kg N over 93 d in 2023---at daily rates of 7.5 and 5.2 g N m$^{-3}$ d$^{-1}$ and efficiencies of 57\% and 64\%, with nightly batching giving 100\% instantaneous removal \parencite{DíazGarcía2025}. Electrical stimulation improved denitrification at 500 mA (carbon anode--cathode best) but not at 100 mA, increasing efficiency at the cost of current efficiency and raising removal cost \parencite{Law2018}. Heated tourmaline both reduced clogging and gave the highest removal across tested HRTs \parencite{Ma2021}. Scaling up, a three-bed catchment-scale Minnesota bioreactor treated 55\% of catchment discharge and removed 713 kg nitrate-N (12.5\% of catchment load; 23\% of throughflow load), but sediment clogging during high flow required a sediment-exclusion system \parencite{Feyereisen2023}. Beyond tile drainage, spring bioreactors in the Chesapeake region averaged 41\% annual load reduction (up to $>$65\% in summer), with hundreds of springs estimated to remove 322--771 kg d$^{-1}$ at low unit cost \parencite{Easton2019}. Non-biological alternatives include anion-exchange resin, which removed nitrate at both low and high flow with near-identical breakthrough volumes but was hindered by sulfate selectivity and regenerant disposal \parencite{MAGETTE1990}. Reviews add that carbon dosing and woodchip--corncob mixtures can substantially raise performance, though carbon dosing risks unused effluent carbon, bioclogging, or pollution swapping \parencite{Wang2025b}.

\subsection{Pollution swapping and greenhouse-gas by-products}

Nonideal conditions can generate N$_2$O, CH$_4$, CO$_2$, sulfate reduction, and methylmercury, but reported risks are generally low \parencite{Hartfiel2022}. For nitrous oxide, column and field studies consistently find emissions small relative to nitrate removed: N$_2$O production was 0.003--0.028\% of N denitrified in columns \parencite{Greenan2009}; an Iowa wall showed no significant N$_2$O difference from controls, indirect emission factors of $\sim$0.0085--0.0088 kg N$_2$O-N kg$^{-1}$ NO$_3$-N, and indirect losses equivalent to tile drainage without the bioreactor \parencite{Moorman2010}; and a Virginia bioreactor had an average N$_2$O emission factor of 0.07\% (median 0.01\%), well below the 0.75\% IPCC comparison \parencite{Bock2018}. Across five Illinois bioreactors, dissolved N$_2$O was produced at four sites but consumed at one, surface fluxes were much smaller than dissolved fluxes, only a small fraction of removed nitrate was emitted as N$_2$O, and modelled watershed-scale indirect emissions were lower with a bioreactor than without \parencite{Brunton2024}. Methane and sulfate responses are more variable. In Illinois, three of five bioreactors produced notable CH$_4$ (driving global-warming potential where present) while soil covers significantly lowered surface CH$_4$ \parencite{Brunton2024}; in an Iowa dual-chamber system both chambers were significant CH$_4$ sources and only one removed sulfate and dissolved N$_2$O, yet the landscape-scale impact was a $<$1\% increase in CO$_2$e, with pollution swapping more likely above 50-h HRT \parencite{Hartfiel2023b}. Sulfate reduction under non-N-limiting conditions is common and increases with HRT---occurring in 83\% of Illinois runs \parencite{Bell2015}, highest in activated-carbon media and associated with lower nitrate removal at long HRT \parencite{Povilaitis2020}, and indicating strongly reduced, potentially H$_2$S-producing conditions in New Zealand \parencite{Rivas2020}. Carbon leaching, especially at start-up, is a further concern: DOC shows an initial flush then decline \parencite{Abusallout2017,Rivas2020}, is dominated by hydrophobic, high-molecular-weight (10--100 kDa) fractions that are strong disinfection-by-product precursors with TOX yields comparable to Suwannee River Fulvic Acid, and can be removed by steel-chip filters or coagulation (alum and PAC more effective than ferric chloride) \parencite{Abusallout2017}. Ammonium and DOC also leach at higher flow \parencite{Greenan2009}, and a catchment-scale system net-exported 212 kg C (22\% of inlet DOC load) \parencite{Feyereisen2023}.

\subsection{Reactive media for phosphorus removal: classes and sorption capacity}

Because woodchips remove little phosphorus---and often leach it at start-up---phosphorus co-treatment has driven the development of coupled and dedicated reactive media \parencite{Christianson2021}. DRP is difficult to intercept because it moves in solution rather than attached to sediment, and reactive filter media offer an edge-of-field means of retaining it. The evidence converges on a two-family conceptual model: iron- (Fe-)based media (iron-oxide-coated sand, iron oxyhydroxides, zero-valent iron, iron filings, steel by-products) retain phosphate primarily by strong inner-sphere adsorption, whereas calcium- (Ca-)rich media (blast-furnace and steel slag, limestone, lime-amended pellets, seashells) retain it primarily by Ca-phosphate precipitation driven by dissolved $\mathrm{Ca^{2+}}$ and elevated pH. Across systematic comparisons, sorption is the preferred DRP-retention mechanism at the edge of tile-drained catchments, and Fe-based materials are generally more suitable than Ca-based because they retain phosphate more stably and over a wider pH range \parencite{Mendes2022,Li2020}.

\subsubsection{Iron-based media}

Iron-based materials are the best-performing and most consistently favoured class. In replicated laboratory comparisons of four iron-based filter materials, each removed at least 95\% of the soluble phosphorus originally present at 1, 10, and 100 ppm $\mathrm{PO_4^{3-}}$-P; at 1000 ppm the ranking was iron oxyhydroxide (IOH) $>$ sand--iron mixture (SMI) $>$ zero-valent iron (ZVI) $>$ powdered iron composite (PIC), and IOH held a slight edge for drainage-water treatment \parencite{Allred2014}. Earlier batch work showed at least 94\% phosphate removal for all ZVI and SMI materials, and columns confirmed substantial removal by SMI3 and ZVI1 even at moderate-to-high concentrations and very short exposure times \parencite{Allred2012}. Synthetic goethite iron oxyhydroxide (SG-IOH) almost totally removed phosphate ($>$98\%) from spiked drainage in saturated columns, with the original formulation outperforming the regenerated one \parencite{Allred2017}. Among adsorbents compared at identical conditions, iron filings reached the highest capacity of 4.95 mg/g and steel slag 1.68 mg/g, versus only 0.10--0.13 mg/g for three natural minerals at 20\,$^{\circ}$C---an order-of-magnitude advantage for industrial by-products \parencite{Sellner2019b}. In long-term columns, the Fe-based CFH material retained the highest P (99\% of high inlet P, 98\% of low) and was most robust under variable flow and loading, whereas calcined diatomite earth performed moderately but consistently \parencite{Pugliese2023}. Granular iron-oxide media (GNEM) removed soluble reactive phosphorus (SRP) below detection far longer than steel by-product media in continuous columns, reaching a final saturated capacity of 1.176 mg P/g \parencite{Hauda2020}. Reviews reinforce this ordering---Fe-based materials generally achieved DRP retention $\geq$95\% in batch and $>$90\% in most column tests, with negligible desorption and strong, stable bonds \parencite{Mendes2022}. Iron-activated alumina (Alcan) achieved up to 94.4\% phosphate removal (maximum capacity 2.608 mg/g) and was proposed as a novel DRP adsorbent for subsurface drainage \parencite{Kumari2025}.

\subsubsection{Calcium-rich slag and lime-based media}

Calcium-rich media are widely available and effective, but performance is governed by the finite supply of reactive $\mathrm{Ca^{2+}}$ and buffering capacity. Granulated blast-furnace slag (GBFS) and cement-kiln-dust (CKD)/silica-sand blends exhibited $>$98\% $\mathrm{PO_4^{3-}}$ removal across batch experiments spanning 3--68 mg/L \parencite{Agrawal2011}. Furnace slag removed all available P at 1--5 mg P/L, with fine particles removing all available P up to 100 mg P/L and following the order FS\_Fine $>$ FS\_SM $>$ FS\_LG at 600 mg P/L; in columns all furnace-slag treatments outperformed sandy-loam-coated pea gravel, with average P removal of 50\% (FS\_SM), 22\% (FS\_LG), and 27\% (FS\_C) \parencite{Sanford2016}. Steel slag reached 1.68 mg/g, and mixtures of small steel chips with limestone yielded 0.95, 1.42, and 2.28 mg/g at 10\%, 20\%, and 40\% steel content \parencite{Sellner2019b}. Fly-ash ceramic pellets amended with lime and bentonite were engineered to raise adsorption while limiting heavy-metal leaching: a lime-free pellet reduced solution P by 80\% within 30 s but leached excessive Cr, whereas the highest-lime pellet had the greatest porosity (14\%), saturated capacity (1.98 mg P/g), and lowest Cr leaching (30 $\mu$g/L) \parencite{Li2017}. Reviews report that Ca-based media commonly achieved 100\% DRP retention at low concentrations but declined at higher concentrations, with some materials (furnace slag, Leca) showing column efficiencies too low for the field; Filtralite-P and seashells performed well at high inlet P but poorly at low inlet P, and limestone showed the lowest retention and highest sensitivity to loading \parencite{Mendes2022,Pugliese2023}.

\subsubsection{Natural minerals and other engineered media}

Natural minerals and unmodified sorbents span a wide, generally lower range. Zeolite and silica sand gave statistically similar, modest $\mathrm{PO_4^{3-}}$ removal of 21--58\%, and coconut-shell activated carbon (CS-AC) removed 70--79\% \parencite{Agrawal2011}. Modified clinoptilolite zeolite improved removal by a factor of 13--18 relative to raw samples, and in drainage-model experiments zeolite placed around the drains achieved the highest removals of the tested configurations (29.28\% phosphate, alongside 59.72\% nitrate and 77.47\% salinity) \parencite{Sadeghi2021}. Comprehensive review of inorganic adsorbents found that phosphate adsorption capacity varies widely among classes and is strongly affected by pH, ionic strength, and coexisting ions; some metal-loaded resins reached 2000--3000 mM/kg at near-neutral pH and generally outperformed plain metal oxyhydroxides and ZVI, while ferrihydrite, lepidocrocite, selected Zr/La materials, layered double hydroxides, and hybrid anion-exchange resins retained relatively strong performance near neutral pH \parencite{Li2020}.

\subsection{Phosphorus retention mechanisms and governing conditions}

\subsubsection{Sorption versus precipitation}

The two families retain phosphate by mechanistically distinct routes. For Fe-based media, retention is dominated by inner-sphere adsorption forming strong, stable bonds with negligible desorption \parencite{Mendes2022}. Isotherm and thermodynamic analyses support chemical binding: best-fit isotherms were Dubinin--Radushkevich for GBFS and zeolite, Freundlich for silica sand, Frumkin for CS-AC, and Langmuir for both CKD/sand blends, with mean free-energy values indicating ion-exchange/chemical sorption for GBFS and the 10:90\% CKD/sand blend but physical sorption for zeolite, silica sand, and CS-AC; the separation factor satisfied $0 < R_L < 1$ and decreased with rising initial concentration \parencite{Agrawal2011}. The Langmuir model generally fit steel and mineral media better than Freundlich, and all materials showed increasingly negative Gibbs free energy and positive enthalpy with rising temperature, indicating endothermic adsorption \parencite{Sellner2019b}. Alcan followed a Freundlich isotherm ($R^2=0.991$) and pseudo-second-order kinetics ($R^2=0.99$), supporting chemisorption \parencite{Kumari2025}, and for lime-amended pellets the Langmuir--Freundlich model fit best, with saturated capacity rising from 0.69 to 1.98 mg/g as lime increased \parencite{Li2017}. For Ca-based slag, retention proceeds by calcium phosphate precipitation: downstream slag segments were more efficient than upstream because they carried higher $\mathrm{Ca^{2+}}$ and pH, and phosphorus was removed in a moving front as $\mathrm{Ca^{2+}}$ and pH-buffer capacity are consumed \parencite{Wang2020d}. Consistent with this, lower pH improved DRP sorption for Fe-based media whereas higher/basic pH improved precipitation-based retention for Ca-based media \parencite{Mendes2022}. The pH dependence of Ca systems is sharp in the field: treated-water pH declined over time and dissolved-P removal dropped dramatically once pH fell below 8.5, with 100\% removal sustained only until cumulative loading reached about 500 mg P/kg slag \parencite{Penn2020}. The same buffering principle underlies subsurface CaCO$_3$ precision treatments, which produced an immediate rise in discharge-water pH by targeting hydrologically active macropores and forming a buffering coating near the drains \parencite{Dalhem2019}.

\subsubsection{pH, coexisting ions, temperature, and particle size}

Media performance is strongly modulated by water chemistry and grain size. Across inorganic adsorbents, phosphate adsorption decreased with increasing pH and acidic conditions generally favoured adsorption, although several materials retained strength near neutral pH \parencite{Li2020}. For lime-amended pellets the opposite held---adsorption increased with pH, with equilibrium sorbate rising from 0.3--1.37 mg/g at pH 4 toward theoretical values at pH 10---reflecting the precipitation pathway \parencite{Li2017}. Coexisting sulfate and bicarbonate often reduced phosphate removal, whereas chloride and nitrate usually had little effect for inner-sphere adsorbents \parencite{Li2020}; correspondingly, nitrate and sulfate caused little interference for steel and mineral media (variations typically $<$10\%), but humic acid and woodchip DOC reduced adsorption by 5--47\%, and real drainage waters caused reductions of 2.3--9.7\% (Arlington) to 9--21\% (Baltic) \parencite{Sellner2019b}. Temperature has a measurable positive effect: adsorption capacities increased by factors of 1.2--2.5 as temperature rose from 5 to 30\,$^{\circ}$C, consistent with the endothermic process \parencite{Sellner2019b}. Smaller particle sizes generally increased DRP sorption and retention and could prolong media lifetime \parencite{Mendes2022}, a result echoed by findings that finer particle size and the presence of Fe oxides enhanced sorption \parencite{Pugliese2023} and that furnace-slag fine particles achieved the greatest sorption \parencite{Sanford2016}. Successful field use therefore requires attention to near-neutral pH performance, selectivity, stability, reusability, recovery, and cost \parencite{Li2020}.

\subsection{Kinetics, longevity, and field-scale phosphorus removal}

\subsubsection{Kinetics and hydraulic behaviour}

Phosphate uptake by the better media is rapid, permitting short contact times. Contact times of only 10--15 s sufficed for near-complete removal from drainage spiked with 1 ppm phosphate-P \parencite{Allred2017}, and iron-based materials reduced soluble P below detection or removed $\geq$90\% even at high flow and short contact times \parencite{Allred2014}. Limestone adsorbed phosphate very rapidly, reaching 98\% of its 24-h capacity within 10 min, while steel slag, small steel chips, and medium steel chips reached 73\%, 67\%, and 51\% of their 24-h capacities within 6 h \parencite{Sellner2019b}. Lime-free fly-ash pellets reduced solution P by 80\% in 30 s, whereas lime-bearing pellets required 15--28 min \parencite{Li2017}. Both granular iron-oxide media and steel-furnace slag removed SRP below detection for up to five hours; after 24 h the slag removed about 50\% while the iron-oxide media still removed below detection, and after continuous operation the slag saturated within seven days whereas the iron-oxide media still removed 20\% SRP after ten days \parencite{Hauda2020}. Removal generally increased with residence and empty-bed contact time (EBCT): in slag columns the difference between 50 s and 17 s was significant but that between 17 s and 7 s was not \parencite{Sanford2016}; longer EBCT improved removal, extended exhaustion time, and raised modelled capacity for Alcan \parencite{Kumari2025}; and P removal rose with retention time in staged slag segments \parencite{Wang2020d}. However, granular iron-oxide capacity was minimally affected by flow rate and by influent SRP typical of subsurface drainage \parencite{Hauda2020}. Adequate hydraulic conductivity underpins field use: all four iron-based materials exceeded the stormwater sand-filter criterion of $1.0\times10^{-3}$ cm/s \parencite{Allred2014}, ZVI and SMI generally exceeded it \parencite{Allred2012}, and silica sand, blast-furnace slag, zeolite, CS-AC, and the CKD/sand blends fell in or near the target range \parencite{Agrawal2011}.

\subsubsection{Desorption, regeneration, and longevity}

The durability of retention differs sharply between families. Fe-based media exhibit low desorption: phosphate on the four iron-based materials was not readily released, especially for SMI and IOH \parencite{Allred2014}, and steel chips showed no apparent desorption while limestone released 9.6\% and steel slag only 0.6\% over ten days \parencite{Sellner2019b}. Reviews note that Fe-based media show negligible desorption whereas Ca-based media more often exhibit weakly bound DRP prone to release, particularly at low DRP or high flow \parencite{Mendes2022}. Several media are regenerable: a NaOH flush plus water rinse remobilised about 81\% of adsorbed phosphate from SG-IOH \parencite{Allred2017}, 0.1 M KOH restored Alcan capacity above 93\% over six cycles \parencite{Kumari2025}, and no desorption was detected from lime-amended pellets after 24 h \parencite{Li2017}. Spent field slag retained useful capacity---with a 10-min retention time it still had considerable P-sorption potential, and laboratory tests showed little water-extractable P release (seven of nine samples released none) \parencite{Shedekar2020}. Longevity is ultimately set by cumulative loading and the moving-front consumption of $\mathrm{Ca^{2+}}$ and buffering capacity in slag \parencite{Wang2020d}: dissolved-P removal remained at 100\% only until roughly 500 mg P/kg slag before becoming variable and sometimes negative \parencite{Penn2020}. Field structures also underperform their design basis: at equal cumulative loading of 200 mg/kg a field structure removed 50 mg/kg (30\%) of dissolved P versus 125 mg/kg ($\sim$65\%) for the pre-construction slag used in design \parencite{Shedekar2020}.

\subsubsection{Field-scale removal structures and drainage models}

Field structures demonstrate that reactive media can treat real drainage discharge while confirming the retention-time and pH limits found in the laboratory. A subsurface tile-drain P removal structure removed 19.2 kg of dissolved P (55\% of cumulative input), while a companion ditch structure removed 0.9 kg (37\% of inflow); precipitation-induced events delivered nearly 97\% of the load while accounting for 65\% of flow volume, and both underperformed laboratory flow-through experiments, with the ditch structure showing clogging and non-uniform flow \parencite{Penn2020}. A ditch structure using aluminium-treated steel slag removed about 27\% of dissolved P and 50\% of total P (2.4 and 9.4 kg), lowering outlet flow-weighted means by 32.2\% (dissolved P) and about 41\% (total P); removal was strong early but highly variable after about January 2017, and the Al-treated slag outperformed non-treated slag at the same site \parencite{Shedekar2020}. A relatively small structure containing 1700 kg of activated Al, operated bottom-up and saturated, treated the entire flow without restricting drainage and removed 30\% of the 2-yr dissolved-P load (240 g) and 29\% of the cumulative nitrate-N load (8.8 kg); during a 9-day dissolved-P spike test it removed 74\% of dissolved P (770 g), about three times the mass removed during the natural-load period \parencite{Penn2024}. In the initial field test of SG-IOH-O, phosphate removal averaged 89\% over the first 200 days before declining to 40\% over the last 80 days, and averaged 75\% versus 34\% for the regenerated material, with no adverse effect on water pH \parencite{Allred2017}. Not all trials succeeded: a GBFS/CS-AC/CKD/sand/zeolite filter blend in a Texas field did not significantly reduce $\mathrm{PO_4^{3-}}$ and instead offloaded phosphate above flow rates of 0.037 L s$^{-1}$, although the same blend removed chlorothalonil (median decrease 69\%, up to 96\%) \parencite{Agrawal2011b}. Flow-through and drainage-model studies provide the design basis: staged slag experiments showed discrete removal declining with time and loading (single-segment removal averaging 33\% at 2.5 mg L$^{-1}$ and 24\% at 5.0 mg L$^{-1}$), rising with number of segments and retention time, and yielded an estimated slag capacity of about 60--61 mg kg$^{-1}$ at 30-min retention time, used to size a filter unit \parencite{Wang2020d}. A continuous drainage model treating actual farm drainage confirmed that clinoptilolite around the drains gave the best performance, with \textit{Phragmites australis} raising removal by a further 10--20\% \parencite{Sadeghi2021}.

\subsection{Coupled treatment, co-benefits, and cost-effectiveness}

\subsubsection{Phosphorus co-removal in bioreactors and stacked designs}

Coupling P-sorbing media with woodchips addresses the woodchip P deficit. Steel-based media are consistently effective P sorbents: steel turnings achieved $>$90\% orthophosphate removal, and placing iron both upstream and downstream of woodchips overcame arrangement conflicts to enhance both nitrate and orthophosphate removal (woodchips reduce nitrate best downstream of iron, while iron re-captures P released by woodchips) \parencite{Goodwin2015}. Steel byproducts showed batch Langmuir $q_{max}$ of 1.30--5.81 mg P/g and a continuous-flow capacity of 3.70 mg P/g, completely removing residual phosphate and reducing effluent nitrite in a two-stage woodchip-then-steel system \parencite{Hua2016}. Fly-ash pellets removed on average 71.4\% of orthophosphate while paired woodchip--pellet configurations held total-phosphate removal at 81.7--84.5\% \parencite{Li2018b}, and iron filings cut drainage P loads by 19\% alongside nitrate removal \parencite{Oladeji2023}. Field bioreactors can become net P sinks over time: Quebec woodchips-only and mixed-media systems released P initially but became sinks within a year, with mixed media removing $\sim$19 times more total-P load and cutting soluble reactive P by 97\% (versus 53\% for woodchips only), though neither met the 0.030 mg L$^{-1}$ criterion \parencite{Husk2018}. A New Zealand bioreactor unexpectedly retained DRP in its second year \parencite{Rivas2020}, and a catchment-scale system removed DRP and TP for most events but exported TP at peak flow \parencite{Feyereisen2023}. Some systems are inconsistent: a Minnesota modular design acted as both P sink and source depending on month and media (fragmented limestone $>$ crushed concrete $>$ steel slag for apparent P-sorbing capacity), with media requiring annual replacement \parencite{Strock2025}, and P removal was not statistically significant in the Virginia system \parencite{Bock2018}. A promising biological route is drying--rewetting: cyclic draining and re-flooding of woodchip columns more than doubled orthophosphate removal (peaking 10--12 h after re-flooding), driven by biological uptake---confirmed because autoclaved woodchips did not take up phosphate and \textit{ppk} transcripts of \textit{Accumulibacter} and \textit{Pseudomonas} peaked with uptake \parencite{Zhang2024b}, with microbial uptake also explaining variable DRP behaviour in columns \parencite{Soupir2018}.

The stacked bioreactor/P-removal design illustrates the hydraulic trade-off: a stacked woodchip bioreactor removed 93\% of the nitrate-N and 29\% of the dissolved P entering it, but because 93\% of the 21-ML total tile discharge bypassed the bioreactor this translated to only a 5.4\% nitrate-N load reduction and 15 g of dissolved P removed, whereas the P-removal structure treating the full flow removed nearly ten times more nitrate-N \parencite{Penn2024}. Reactive media can also address multiple contaminants: in batch tests three SMI materials removed at least 86\% of nitrate and ZVI1 achieved an 87.8\% atrazine decrease, and a two-stage SMI3-then-ZVI1 column removed nitrate, phosphate, and atrazine together \parencite{Allred2012}. Because Fe-based P adsorbents produce negligible acidity or alkalinity effects, they pair naturally with woodchip bioreactors, and steel by-product filters were proposed for such coupling \parencite{Sellner2019b}. Complementary edge-of-field practices target nitrate directly: all six saturated riparian buffers effectively removed nitrate from tile-outlet water, with annual loads removed of 13--179 kg N over drainage areas of 3.4--40.5 ha, removal effectiveness of 8--84\%, and rates of 0.27--4.13 g N m$^{-1}$ d$^{-1}$ \parencite{Jaynes2019}.

\subsubsection{Secondary benefits and monitoring context}

Beyond nutrients, these systems provide co-benefits and rest on well-defined monitoring. Woodchip columns effectively removed \textit{Salmonella} and \textit{E.\ coli} under all conditions (greater at 21.5 than 10\,$^{\circ}$C), indicating a pathogen-removal co-benefit \parencite{Soupir2018}. Reliable assessment depends on discrete influent and effluent points \parencite{Easton2019} and on standardised monitoring set-ups and minimum reporting measures \parencite{Christianson2021}. Supporting analytical methods have been validated---the Dumas method for total N and $^{15}$N of ground plant material and concentrated drainage (highest recovery 97.4\%, improved with copper shot) \parencite{Schindler1999}---and resin-capsule methods can intercept nitrate moving via macropore and bypass flow, the primary transport pathway to subsurface drainage \parencite{Pampolino2000}. Related drainage-chemistry work situates reactive-media P removal within broader management: subsurface CaCO$_3$ precision treatments in acid sulfate soils kept pH $+0.4$ units above reference four years later, decreased aluminium by about 95\%, and shifted Fe speciation toward schwertmannite and amorphous oxides, whereas Ca(OH)$_2$ treatments cut Al by up to 99\% immediately but were short-lived \parencite{Dalhem2019}. CaCO$_3$ amendment also helped reduce swelling and excess field-capacity water retention \parencite{He2015}. Where drainage doubles as subirrigation supply, a new performance method reported distribution uniformity of the low quarter of 48--78\% that improved with longer durations or slower flow \parencite{Jia2017}. Subsurface pipe drainage markedly altered soil ion contents (lowering $\mathrm{Na^+}$, $\mathrm{Ca^{2+}}$, $\mathrm{Mg^{2+}}$ while $\mathrm{SO_4^{2-}}$ became most abundant), with HYDRUS-2D/UnsatChem reproducing ion discharge ($R^2 > 0.87$) and showing ion-composition-based designs outperform total-salt designs \parencite{Liu2025b}. For monitoring, field electrical conductivity was a precise, cost-effective estimator of major-ion concentration, with $\mathrm{HCO_3^-}$, $\mathrm{Ca^{2+}}$, and $\mathrm{Mg^{2+}}$ accounting for 90\% of calculated conductivity \parencite{Zimmerman2017}, and electrochemical treatment of drainage brine was optimised by response-surface methods to 3.6 mol L$^{-1}$ brine, pH 9, 12.2 V, and 22.5-cm electrode depth \parencite{Yang2023}.

\subsubsection{Cost-effectiveness and comparison with other practices}

Bioreactors are generally low-cost per unit nitrogen removed. Installation costs range from under \$5{,}000 to \$27{,}000, with cost efficiencies generally from $<$\$2.50 to $\sim$\$20 kg$^{-1}$ N y$^{-1}$ (up to \$48) \parencite{Christianson2021}. Corncob (CC75) systems were the most cost-efficient carbon option (\$10.56--13.89 kg$^{-1}$ N versus up to \$88.11 for woodchips) \parencite{Law2023}. For pumped systems, unit cost fell as lifespan, mass-removal rate, and operating period rose---raising mass-removal rate from low to high cut unit cost by 91\%---with best-case costs of $\sim$\$4.80--23.81 kg$^{-1}$ NO$_3$-N \parencite{Hartfiel2024}. Spring bioreactors achieved \$0.54--7.60 kg$^{-1}$ N y$^{-1}$ (typically $<$\$5 for higher-flow, higher-concentration springs) \parencite{Easton2019}, whereas electrical stimulation raised costs above the base case \parencite{Law2018}. For phosphorus, reactive-media treatment is feasible and optimisable by media choice and end-of-life strategy: for a subsurface case study, the media volume to remove 1.66 kg SRP was 4.22 m$^3$ for granular iron-oxide media versus 15.5 m$^3$ for steel-furnace slag, with annual cost per hectare of \$1641 (scenario A) to \$906 (scenario C), the lowest being single-use slag disposed after use \parencite{Hauda2020}. Recycled steel by-products were identified as highly efficient phosphate adsorbents across wide temperature, pH, organic-matter, and drainage-water conditions and recommended for subsurface drainage \parencite{Sellner2019b}; furnace-slag fine particles may suit an inline system \parencite{Sanford2016}; durable fly-ash pellets could serve as a tile-outlet filter \parencite{Li2017}; Fe-rich materials provide extended, relatively high P retention with lower sensitivity to variable flow and concentration \parencite{Pugliese2023}; and readily acquired slag/zeolite constituents remain attractive for large-scale use despite offloading observed for one blend \parencite{Agrawal2011b}. Saturated riparian buffers had a mean equalised annual cost of US\$213.83 (US\$2.94 kg$^{-1}$ N removed) and are simple, cheap, and quick to install \parencite{Jaynes2019}.

Set against other edge-of-field practices, bioreactors are effective for nitrate but comparatively limited in co-benefits. A US Corn Belt review of 119 records found all five practices (controlled drainage, water-quality wetlands, bioreactors, drainage water recycling, saturated buffers) effective at removing nitrate but with differing additional benefits: water-quality wetlands gave the highest nitrate reduction plus phosphorus removal and habitat value, drainage water recycling offered the greatest yield and co-benefit potential, and denitrifying bioreactors removed nitrate effectively but with limited extra benefits and a need for further study of phosphorus, pesticides, and greenhouse gases \parencite{Mitchell2023}. Global synthesis places bioreactor flow-season removal at 14--66\%, alongside constructed wetlands ($\sim$39\% average TN removal) and saturated buffers and controlled drainage ($\sim$30--50\% annual N-load reduction), while noting that agricultural drainage's high-nitrate, low-carbon, low-temperature character complicates transfer of municipal treatment approaches \parencite{Wang2025b}. Constructed-wetland and paired-system studies reinforce that design and sizing dominate outcomes: about 40--50\% N removal is generally achievable at area ratios of 1--4\% with minor gains beyond \parencite{Mendes2021}, and paired surface-flow wetland plus woodchip systems acted as nutrient sinks (TN removal 8--51\%, TP 0--80\%), with the woodchip component contributing more removal than the sedimentation basin and requiring smaller footprints than standalone wetlands \parencite{Pugliese2020}. Collectively, this evidence establishes woodchip bioreactors and reactive-media P filters as dependable, cost-effective nutrient-reduction practices whose performance is well predicted by residence time, temperature, flow, pH, and cumulative loading, and whose principal remaining design frontiers concern phosphorus co-treatment, greenhouse gases, and pesticides.

\printbibliography
\end{document}
